\section{Analysis Strategy Overview}
\label{sec:StrategyOverview}


% master equation 
The analysis strategy centers around measuring the signal branching fractions with respect to the branching fraction of \BpmKpmJpsi, hereon referred to as the reference process.  
Yields of the signal and references processes are measured separately from fits to data in channels consisting of selections designed to yield a high purity in their respective target process. 
These yields are used to compute the signal branching fractions using equation~\ref{eq:BRFormula} 
\begin{eqnarray}
  \hbox{\hspace{-1cm}} \BR (\Bmumu)=&
  \hbox{\hspace{-3mm}}\BR (\BpmKpmJpsi\! \to \mumu K^\pm)\times
  \frac{f_{u}}{f_{s}}\times N_{\mumu}\times \left( N_{J/\psi K^\pm} \theInvRhok \right)^{-1}\,
  \label{eq:BRFormula}
\end{eqnarray}
where 
\begin{itemize}
  \item $N_{\mumu}$ and $N_{J/\psi K^\pm}$ are the yields of the signal and reference processes in the signal and reference channels respectively
  \item $(A\varepsilon)^k_{\mu^+\mu^-}$ and $(A\varepsilon)^k_{J/\psi K^\pm}$ are the products of acceptance and efficiency (defined in Section~\ref{sec:AccepEffRatio}) for the signal and reference processes in their respective channels. 
  \item $f_{u}$ and $f_{s}$ are the production fractions of \Bp and \Bs mesons respectively.
\end{itemize}
The acceptance and efficiencies are computed using Monte Carlo (MC) simulation, while $\BR (\BpmKpmJpsi)$ and $f_{u}$ and $f_{s}$ are taken as external inputs. 
\textred{Specifically, they are taken from...} 
The main benefit of this approach is that the branching fraction measurement is independent of the $p p \rightarrow b\bar{b}$ cross-section which has a large theoretical uncertainty. 
There is also a significant cancelling of many systematic uncertainties common between the signal and reference channels. 

Triggers optimized for the signal and reference processes are used to select events for the signal and reference channels. 
The remaining selections defining the signal channel consist of cuts on the kinematics of muon and $B$-meon candidates. 
This includes requiring the invariant mass of the muon candidates to be fall within $100\,\MeV$ of the mass of the \Bs meson. 
Upper and lower sidebands are defined by events failing this tight $200\,\MeV$  mass window, referred to as the fitting region, but passing a looser $1200\,\MeV$ window.
Data events falling in the sideband regions are used in designing the analysis while data in the tight mass window is blinded until the analysis is approved.  
To reduce what would otherwise be an overwhelmingly dominant combinatorial background, arising from events in which two muons originate from different \bquark decays, a Boosted Decision Tree (BDT) is trained on the signal MC and the data in the sidebands. 
The signal channel fit is then performed simultaneously in bins defined using the BDT score discriminant. 
The selections for the reference channel are designed to closely match those used in the signal channel.  
Studies using toy MC experiments are used to validate the signal and reference channel fits and evaluate systematic uncertainties. 